% ============================================================================
% 03_optimization/05_peft.tex
% Section-body file: include immediately after the corresponding \section
% command in the master document. No \chapter or \section command is defined
% here, and the global bibliography is printed by the master document.
% ============================================================================

The preceding sections have established the optimization procedures and
training objectives used to construct and improve language models. They have
also considered techniques such as quantization and distillation that reduce
the computational or representational cost of training and deployment. A
closely related problem arises after pretraining: a model that has acquired
general-purpose linguistic representations must often be adapted to a
particular domain, task, or application. The central question is whether this
adaptation requires changing every parameter of the pretrained model.

Full fine-tuning provides a direct answer by allowing all model parameters to
change. For modern large language models, however, the resulting memory,
computation, and storage requirements can be substantial. Moreover, when one
pretrained model must serve many downstream tasks, maintaining a complete
fine-tuned copy for every task duplicates most of the original model. Parameter-
efficient fine-tuning (PEFT) addresses this problem by restricting the
task-specific optimization to a small subset of existing parameters, by
introducing lightweight trainable modules, or by reparameterizing the update
so that it is represented with fewer degrees of freedom. The objective is not
to replace the pretrained model with a smaller network, but to adapt its
existing computational structure with a smaller task-specific optimization
problem \parencite{prottasha2025peft,he2025smt}.

\subsection{From Full Fine-Tuning to Parameter-Efficient Adaptation}
\label{subsec:peft-motivation}

A language model can be represented as a parameterized function
\begin{equation}
    f_{\boldsymbol{\theta}}:\mathcal{X}\rightarrow\mathcal{Y},
    \qquad
    \boldsymbol{\theta}\in\mathbb{R}^{P},
    \label{eq:peft-model-function}
\end{equation}
where $\mathcal{X}$ and $\mathcal{Y}$ denote the input and output spaces,
$\boldsymbol{\theta}$ is the vector containing the model parameters, and $P$
is the number of scalar parameters. In a large language model, $P$ may reach
billions. The pretrained parameter vector is denoted by
$\boldsymbol{\theta}_0$. It is obtained by optimizing a pretraining objective
over a large corpus. For an autoregressive language model, a representative
objective for a token sequence $x_1,\ldots,x_T$ is
\begin{equation}
    \mathcal{L}_{\mathrm{pretrain}}(\boldsymbol{\theta})
    =
    -\sum_{t=1}^{T}
    \log p_{\boldsymbol{\theta}}(x_t\mid x_{<t}),
    \label{eq:peft-pretraining-objective}
\end{equation}
where $x_{<t}$ denotes the tokens preceding position $t$. The resulting
$\boldsymbol{\theta}_0$ provides an initialization that already encodes
statistical regularities learned from the pretraining distribution.

For a downstream task, let
\begin{equation}
    \mathcal{D}_{\mathrm{task}}
    =
    \{(x_i,y_i)\}_{i=1}^{N}
    \label{eq:peft-downstream-dataset}
\end{equation}
be a dataset of $N$ input-target pairs. The task-specific empirical risk is
\begin{equation}
    \mathcal{L}_{\mathrm{task}}(\boldsymbol{\theta})
    =
    \frac{1}{N}\sum_{i=1}^{N}
    \ell\!\left(f_{\boldsymbol{\theta}}(x_i),y_i\right),
    \label{eq:peft-task-loss}
\end{equation}
where $\ell$ is the loss appropriate to the task. For classification it may
be cross-entropy; for autoregressive generation it may be token-level
negative log-likelihood. Full fine-tuning solves
\begin{equation}
    \boldsymbol{\theta}^{*}
    =
    \arg\min_{\boldsymbol{\theta}\in\mathbb{R}^{P}}
    \mathcal{L}_{\mathrm{task}}(\boldsymbol{\theta}),
    \qquad
    \boldsymbol{\theta}^{(0)}=\boldsymbol{\theta}_0.
    \label{eq:peft-full-ft}
\end{equation}
The complete pretrained parameter vector is therefore the optimization
variable. At iteration $t$, a gradient method evaluates
\begin{equation}
    \mathbf{g}^{(t)}
    =
    \nabla_{\boldsymbol{\theta}}
    \mathcal{L}_{\mathrm{task}}
    \bigl(\boldsymbol{\theta}^{(t)}\bigr)
    \label{eq:peft-full-gradient}
\end{equation}
and, in the case of gradient descent, updates
\begin{equation}
    \boldsymbol{\theta}^{(t+1)}
    =
    \boldsymbol{\theta}^{(t)}-\eta\mathbf{g}^{(t)},
    \label{eq:peft-gd-update}
\end{equation}
where $\eta>0$ is the learning rate. In full fine-tuning, every component of
$\boldsymbol{\theta}$ is eligible for an update.

The memory cost is not limited to the model weights. Training also requires
gradient storage, optimizer state, and intermediate activations needed by
backpropagation. For Adam, the first and second moment estimates are
\begin{align}
    \mathbf{m}_t
    &=
    \beta_1\mathbf{m}_{t-1}
    +(1-\beta_1)\mathbf{g}_t,
    \label{eq:peft-adam-m}\\
    \mathbf{v}_t
    &=
    \beta_2\mathbf{v}_{t-1}
    +(1-\beta_2)\mathbf{g}_t^{\odot 2},
    \label{eq:peft-adam-v}
\end{align}
where $\beta_1,\beta_2\in[0,1)$ are decay coefficients and $\odot$ denotes
element-wise multiplication. The bias-corrected estimates are
\begin{equation}
    \widehat{\mathbf{m}}_t
    =
    \frac{\mathbf{m}_t}{1-\beta_1^t},
    \qquad
    \widehat{\mathbf{v}}_t
    =
    \frac{\mathbf{v}_t}{1-\beta_2^t},
    \label{eq:peft-adam-bias-correction}
\end{equation}
and the Adam update is
\begin{equation}
    \boldsymbol{\theta}_{t+1}
    =
    \boldsymbol{\theta}_t
    -
    \eta
    \frac{\widehat{\mathbf{m}}_t}
    {\sqrt{\widehat{\mathbf{v}}_t}+\epsilon},
    \label{eq:peft-adam-update}
\end{equation}
where $\epsilon$ is a small positive constant for numerical stability.
Consequently, the optimizer maintains additional quantities whose sizes
scale with the number of trainable parameters.

A useful first-order memory accounting model is
\begin{equation}
    M_{\mathrm{FT}}
    \approx
    P(b_w+b_g+b_m+b_v)+M_{\mathrm{act}},
    \label{eq:peft-full-memory}
\end{equation}
where $b_w$, $b_g$, $b_m$, and $b_v$ are the bytes per parameter used for
weights, gradients, first moments, and second moments, respectively, and
$M_{\mathrm{act}}$ is the activation memory. This is an accounting model, not
a universal implementation-specific formula: precision, optimizer
implementation, activation checkpointing, gradient accumulation, and
offloading change the actual footprint.

The scale of the problem is illustrated by the LLaMA-7B example reported by
\textcite{he2025smt}. In their stated configuration, full fine-tuning without
CPU offloading requires approximately $58$ GB of GPU memory: about $13.6$ GB
for trainable parameters, $40$ GB for Adam optimizer states and gradients, and
$2$ GB for activations. The exact values depend on the implementation and
precision, but the example shows why weight storage alone is an inadequate
estimate of training memory.

Full fine-tuning also has a multi-task storage cost. If a pretrained model
with $P$ parameters is independently fine-tuned for $K$ tasks, storing all
task-specific copies requires approximately
\begin{equation}
    S_{\mathrm{FT}}=KP.
    \label{eq:peft-ft-storage}
\end{equation}
If the base model is shared and each task stores only $q$ adaptation
parameters, where $q\ll P$, the corresponding storage is
\begin{equation}
    S_{\mathrm{PEFT}}=P+Kq.
    \label{eq:peft-storage}
\end{equation}
The ratio is
\begin{equation}
    \frac{S_{\mathrm{PEFT}}}{S_{\mathrm{FT}}}
    =
    \frac{1}{K}+\frac{q}{P}.
    \label{eq:peft-storage-ratio}
\end{equation}
Thus, the storage advantage becomes increasingly important as the number of
downstream tasks grows.

The statistical motivation is related but distinct. A model with a very large
number of adjustable parameters can fit complex patterns, but a downstream
dataset may contain relatively few examples. In such a setting, unrestricted
adaptation can increase estimation variance and overfitting. The PEFT survey
describes restricted, structured updates as a form of implicit
regularization that can be beneficial in low-resource settings
\parencite{prottasha2025peft}. Restricting updates can also reduce interference
with pretrained knowledge and thereby mitigate, but does not guarantee the
elimination of, catastrophic forgetting. The central design question is
therefore how to restrict the update space without preventing the model from
learning the task-specific information it needs.

\paragraph{General formulation and taxonomy.}

The common mathematical structure of PEFT is to replace unrestricted
optimization over $\boldsymbol{\theta}\in\mathbb{R}^{P}$ with optimization
over a smaller task-specific parameter vector
$\boldsymbol{\phi}\in\mathbb{R}^{q}$, where $q\ll P$. A general update
parameterization is
\begin{equation}
    \Delta\boldsymbol{\theta}
    =
    g(\boldsymbol{\phi}),
    \qquad
    \boldsymbol{\theta}_{\mathrm{eff}}
    =
    \boldsymbol{\theta}_0+g(\boldsymbol{\phi}),
    \label{eq:peft-general-update}
\end{equation}
where $g$ maps the lower-dimensional adaptation variables into an update of
the pretrained parameters. The optimization problem becomes
\begin{equation}
    \boldsymbol{\phi}^{*}
    =
    \arg\min_{\boldsymbol{\phi}\in\mathbb{R}^{q}}
    \frac{1}{N}\sum_{i=1}^{N}
    \ell\!\left(
        f_{\boldsymbol{\theta}_0+g(\boldsymbol{\phi})}(x_i),y_i
    \right).
    \label{eq:peft-reduced-objective}
\end{equation}
The pretrained parameters may remain fixed, while gradients are propagated
through the model to optimize $\boldsymbol{\phi}$. In methods that add modules
rather than explicitly reparameterize the original weights, the model can be
written more directly as
$f_{\boldsymbol{\theta}_0,\boldsymbol{\phi}}(x)$.

Selective fine-tuning can be represented using a binary mask
$\mathbf{M}\in\{0,1\}^{P}$:
\begin{equation}
    \boldsymbol{\theta}_{\mathrm{eff}}
    =
    \boldsymbol{\theta}_0+
    \mathbf{M}\odot\Delta\boldsymbol{\theta},
    \qquad
    q=\|\mathbf{M}\|_0,
    \label{eq:peft-selection-mask}
\end{equation}
where $\|\mathbf{M}\|_0$ counts the selected scalar parameters. A value of
one permits an update, while zero freezes the corresponding parameter.
Other PEFT families do not necessarily use a binary mask, but they share the
principle of reducing the number of independent adaptation variables.

The survey by \textcite{prottasha2025peft} organizes PEFT techniques around
additive, selective, reparameterized, hybrid, and unified approaches, and
separately discusses mixture-of-experts-based variants. These categories
describe different ways of controlling the adaptation space. Additive
methods introduce task-specific modules or representations. Selective methods
choose a subset of existing parameters. Reparameterized methods express
weight changes through lower-dimensional variables. Hybrid and unified
frameworks combine multiple mechanisms, while mixture-of-experts designs
route inputs or tokens among specialized adaptation components. Figure
\ref{fig:peft-taxonomy} summarizes this conceptual organization.

\begin{figure}[htbp]
    \centering
    % Original conceptual taxonomy; uses the standard LaTeX picture environment.
\setlength{\unitlength}{1mm}
\begin{picture}(150,58)
  \put(58,47){\framebox(34,9){\textbf{PEFT}}}
  \put(75,47){\line(0,-1){7}}
  \put(14,40){\line(1,0){120}}
  \put(14,40){\line(0,-1){8}}
  \put(44,40){\line(0,-1){8}}
  \put(74,40){\line(0,-1){8}}
  \put(104,40){\line(0,-1){8}}
  \put(134,40){\line(0,-1){8}}
  \put(0,14){\framebox(28,18){\shortstack{\scriptsize\textbf{Additive}\\\scriptsize adapters\\\scriptsize / prompts}}}
  \put(30,14){\framebox(28,18){\shortstack{\scriptsize\textbf{Selective}\\\scriptsize parameter\\\scriptsize selection}}}
  \put(60,14){\framebox(28,18){\shortstack{\scriptsize\textbf{Reparameterized}\\\scriptsize low-rank\\\scriptsize updates}}}
  \put(90,14){\framebox(28,18){\shortstack{\scriptsize\textbf{Hybrid / unified}\\\scriptsize combined\\\scriptsize mechanisms}}}
  \put(120,14){\framebox(28,18){\shortstack{\scriptsize\textbf{MoE-based}\\\scriptsize routed\\\scriptsize experts}}}
\end{picture}

    \caption{Conceptual taxonomy of parameter-efficient fine-tuning. The
    branches summarize the principal families discussed by
    \textcite{prottasha2025peft}; the diagram is an original schematic rather
    than a reproduction of the survey figure.}
    \label{fig:peft-taxonomy}
\end{figure}
% Alt text for accessibility register: A root node labelled PEFT branches into
% additive, selective, reparameterized, hybrid/unified, and mixture-of-experts
% adaptation families, each with a short description of its mechanism.

The categories should not be interpreted as mutually exclusive mathematical
classes. A method can combine an additive module with a low-rank update, or
use a sparse selection rule to determine where a reparameterized update is
applied. The classification is useful because it identifies the principal
mechanism that controls the trainable parameter space.

\subsection{Additive and Prompt-Based Adaptation}
\label{subsec:peft-additive-prompts}

Additive fine-tuning introduces trainable modules or representations into
the existing model while retaining the pretrained backbone. The added
parameters encode task-specific transformations, and the original model
parameters can remain frozen. Adapters and soft prompts are two important
examples \parencite{prottasha2025peft}.

\paragraph{Serial adapters.}
A serial adapter is commonly implemented as a bottleneck transformation.
Let $\mathbf{h}\in\mathbb{R}^{d}$ be a hidden representation, and let
$d_b\ll d$ be the bottleneck dimension. The down-projection is
\begin{equation}
    \mathbf{h}_{\mathrm{down}}
    =
    \operatorname{ReLU}
    \left(
        \mathbf{W}_{\mathrm{down}}\mathbf{h}
        +\mathbf{b}_{\mathrm{down}}
    \right),
    \qquad
    \mathbf{W}_{\mathrm{down}}\in\mathbb{R}^{d_b\times d},
    \label{eq:peft-adapter-down}
\end{equation}
and the up-projection is
\begin{equation}
    \mathbf{h}_{\mathrm{up}}
    =
    \mathbf{W}_{\mathrm{up}}\mathbf{h}_{\mathrm{down}}
    +\mathbf{b}_{\mathrm{up}},
    \qquad
    \mathbf{W}_{\mathrm{up}}\in\mathbb{R}^{d\times d_b}.
    \label{eq:peft-adapter-up}
\end{equation}
The output is added through a residual connection:
\begin{equation}
    \mathbf{h}_{\mathrm{out}}
    =
    \mathbf{h}+\mathbf{h}_{\mathrm{up}}.
    \label{eq:peft-adapter-residual}
\end{equation}
Ignoring biases, one adapter contains
\begin{equation}
    P_{\mathrm{adapter}}
    =
    dd_b+d_bd
    =
    2dd_b
    \label{eq:peft-adapter-parameters}
\end{equation}
trainable parameters. A dense $d\times d$ transformation contains $d^2$
parameters, so the ratio is $2d_b/d$. If one adapter is inserted in each of
$L$ layers, the approximate count is $2Ldd_b$, subject to the particular
placement and architecture. The bottleneck creates a capacity-efficiency
trade-off: a larger $d_b$ permits a richer transformation but increases
parameters and computation.

\paragraph{Parallel and hybrid adapters.}
A parallel adapter processes a representation alongside the main model
pathway. If its output is
\begin{equation}
    \mathbf{h}_{\mathrm{parallel}}
    =
    \mathbf{W}_{\mathrm{parallel}}\mathbf{h}
    +\mathbf{b}_{\mathrm{parallel}},
    \label{eq:peft-parallel-adapter}
\end{equation}
the combined representation can be written as
\begin{equation}
    \mathbf{h}_{\mathrm{out}}
    =
    \mathbf{h}_{\mathrm{main}}
    +\alpha\mathbf{h}_{\mathrm{parallel}},
    \label{eq:peft-parallel-combine}
\end{equation}
where $\alpha$ controls the contribution of the parallel pathway. Hybrid
adapters combine serial and parallel pathways, for example,
\begin{equation}
    \mathbf{h}_{\mathrm{out}}
    =
    \beta\mathbf{h}_{\mathrm{serial}}
    +\gamma\mathbf{h}_{\mathrm{parallel}},
    \label{eq:peft-hybrid-adapter}
\end{equation}
with coefficients $\beta$ and $\gamma$. The survey discusses serial,
parallel, and hybrid designs, as well as single-task and multi-task adapter
use cases \parencite{prottasha2025peft}. Figure \ref{fig:peft-adapter-prompt}
contrasts the broad idea of inserting a trainable adapter with that of
introducing trainable prompt representations.

\begin{figure}[htbp]
    \centering
    % Original conceptual comparison; uses the standard LaTeX picture environment.
\setlength{\unitlength}{1mm}
\begin{picture}(150,62)
  \put(18,55){\makebox(45,5){\textbf{Adapter-based adaptation}}}
  \put(92,55){\makebox(45,5){\textbf{Prompt-based adaptation}}}

  % Left: frozen backbone with trainable residual adapter.
  \put(3,32){\framebox(24,12){\shortstack{Frozen\\Transformer}}}
  \put(39,32){\framebox(24,12){\shortstack{Frozen\\Transformer}}}
  \put(27,38){\vector(1,0){12}}
  \put(15,32){\line(0,-1){8}}
  \put(7,15){\framebox(16,9){\shortstack{Trainable\\adapter}}}
  \put(15,24){\vector(0,1){8}}
  \put(51,32){\vector(0,-1){8}}
  \put(40,15){\framebox(22,9){\shortstack{Residual\\combination}}}
  \put(23,19){\vector(1,0){17}}
  \put(3,4){\makebox(62,6){\scriptsize Added module modifies hidden states}}

  % Right: learned prompt vectors prepended to token embeddings.
  \put(82,38){\framebox(20,10){\shortstack{Prompt\\vectors}}}
  \put(108,38){\framebox(30,10){\shortstack{Token\\embeddings}}}
  \put(102,43){\vector(1,0){6}}
  \put(95,38){\line(0,-1){10}}
  \put(85,17){\framebox(48,11){\shortstack{Frozen Transformer\\(attention / FFN)}}}
  \put(109,38){\vector(0,-1){10}}
  \put(85,8){\makebox(48,6){\scriptsize Prompt vectors are optimized}}
\end{picture}

    \caption{Conceptual comparison of adapter-based and prompt-based
    adaptation. Adapters add trainable transformations within the model
    pathway, whereas prompt methods add trainable representations to the
    input or attention context. Original schematic based on the mechanisms
    surveyed by \textcite{prottasha2025peft}.}
    \label{fig:peft-adapter-prompt}
\end{figure}
% Alt text for accessibility register: The left pathway shows a frozen
% pretrained backbone with a small trainable adapter residual branch. The
% right pathway shows trainable prompt vectors concatenated with token
% embeddings before entering a frozen Transformer.

\paragraph{Soft prompt tuning.}
Let the token embedding sequence for an input $x$ be
$\mathbf{E}(x)\in\mathbb{R}^{n\times d}$, where $n$ is the number of input
tokens and $d$ is the embedding dimension. Soft prompt tuning introduces
$m$ trainable vectors,
\begin{equation}
    \mathbf{P}\in\mathbb{R}^{m\times d},
    \label{eq:peft-soft-prompt}
\end{equation}
and concatenates them with the input embeddings:
\begin{equation}
    \mathbf{X}_{\mathrm{prompt}}
    =
    [\mathbf{P};\mathbf{E}(x)]
    \in\mathbb{R}^{(m+n)\times d}.
    \label{eq:peft-prompt-concat}
\end{equation}
The pretrained model is frozen and the prompt is optimized:
\begin{equation}
    \mathbf{P}^{*}
    =
    \arg\min_{\mathbf{P}}
    \mathcal{L}
    \left(
        f_{\boldsymbol{\theta}_0}
        ([\mathbf{P};\mathbf{E}(x)]),y
    \right).
    \label{eq:peft-prompt-objective}
\end{equation}
The prompt contains $md$ trainable scalar values. Unlike a discrete prompt,
whose elements are vocabulary tokens, a soft prompt is a continuous matrix
of trainable embeddings and is not required to correspond to ordinary words.

\paragraph{Prefix tuning.}
Prefix tuning extends the prompt principle into the attention computation.
For query, key, and value matrices $\mathbf{Q}$, $\mathbf{K}$, and
$\mathbf{V}$, introduce trainable prefix key and value representations
$\mathbf{P}_K$ and $\mathbf{P}_V$. The attention operation becomes
\begin{equation}
    \operatorname{Attention}
    \left(
        \mathbf{Q},
        [\mathbf{P}_K;\mathbf{K}],
        [\mathbf{P}_V;\mathbf{V}]
    \right)
    =
    \operatorname{softmax}
    \left(
        \frac{
            \mathbf{Q}[\mathbf{P}_K;\mathbf{K}]^{\mathsf T}
        }{\sqrt{d_k}}
    \right)
    [\mathbf{P}_V;\mathbf{V}],
    \label{eq:peft-prefix-attention}
\end{equation}
where $d_k$ is the key dimension. Prompt tuning modifies the input
embedding sequence, whereas prefix tuning makes trainable contextual
representations available within attention. The survey also discusses
multi-layer continuous prompts such as P-Tuning v2, as well as discrete
prompt optimization methods such as RLPrompt and sparse prompt selection
\parencite{prottasha2025peft}.

For a discrete prompt, token identities are not ordinary continuous
optimization variables. One alternative is to optimize a policy over
candidate token sequences. If $\boldsymbol{\psi}$ denotes policy parameters
and $R(P)$ is the reward associated with a prompt $P$, a policy-gradient
update has the form
\begin{equation}
    \boldsymbol{\psi}_{t+1}
    =
    \boldsymbol{\psi}_{t}
    +
    \eta\nabla_{\boldsymbol{\psi}}
    \mathbb{E}[R(P)].
    \label{eq:peft-rlprompt}
\end{equation}
Sparse prompt selection can instead be expressed as
\begin{equation}
    P_{\mathrm{sparse}}^{*}
    =
    \arg\min_{P}\|P\|_0
    \quad
    \text{subject to}
    \quad
    R(P)\geq R_{\mathrm{threshold}},
    \label{eq:peft-sparse-prompt}
\end{equation}
where $\|P\|_0$ counts the selected prompt elements. These formulations
illustrate that prompt-based PEFT can optimize either continuous embeddings
or a policy that generates discrete prompts.

\paragraph{Scaling-based adaptation.}
The survey also describes Propulsion, which modifies feature representations
through trainable scaling values. For feature groups
$\mathbf{V}=[v_1,\ldots,v_s]$, its transformation is expressed as
\begin{equation}
    \mathbf{V}'
    =
    [v_1\odot z_1^k,\,
     v_2\odot z_2^k,\,
     \ldots,\,
     v_s\odot z_s^k],
    \label{eq:peft-propulsion}
\end{equation}
where $z_i$ is a trainable scaling parameter, $k$ is a scaling exponent, and
$\odot$ denotes element-wise multiplication. This mechanism changes the
relative contribution of feature components without introducing a complete
dense transformation.

Additive and prompt-based methods are parameter-efficient, but they do not
eliminate all computation. Adapters add operations to the forward path.
Prompt and prefix methods increase the effective context length; if the
attention sequence grows from $n$ to $n+m$, the dense attention score matrix
grows from $n\times n$ to $(n+m)\times(n+m)$. Its leading pairwise
interaction cost therefore changes from $\mathcal{O}(n^2)$ to
$\mathcal{O}((n+m)^2)$. A small trainable parameter count should consequently
not be confused with a proportionate reduction in inference computation.

\subsection{Selective and Reparameterized Adaptation}
\label{subsec:peft-selective-reparameterized}

Selective fine-tuning does not introduce a new transformation as its defining
operation. Instead, it chooses which existing parameters may change. Partition
the pretrained parameters into selected and frozen components,
\begin{equation}
    \boldsymbol{\theta}
    =
    \{\boldsymbol{\theta}_{s},
      \boldsymbol{\theta}_{f}\},
    \qquad
    \boldsymbol{\theta}_{f}=\text{constant}.
    \label{eq:peft-selected-partition}
\end{equation}
The optimization is
\begin{equation}
    \boldsymbol{\theta}_{s}^{*}
    =
    \arg\min_{\boldsymbol{\theta}_{s}}
    \mathcal{L}
    \left(
        f_{\boldsymbol{\theta}_{s},\boldsymbol{\theta}_{f}}(x),y
    \right).
    \label{eq:peft-selective-objective}
\end{equation}
The key problem is the selection rule.

\paragraph{Gradient and Fisher-information criteria.}
For parameter $\theta_i$, the local sensitivity of the loss is described by
\begin{equation}
    g_i
    =
    \frac{\partial\mathcal{L}}{\partial\theta_i}.
    \label{eq:peft-parameter-gradient}
\end{equation}
A simple importance score is $C(\theta_i)=|g_i|$. Parameters with large
scores may be selected, subject to a trainable-parameter budget. This score
is local: it depends on the current model, data batch, and loss, and is not a
task-independent measure of universal parameter importance.

A statistical alternative is Fisher information. For a model distribution
$p(x;\boldsymbol{\theta})$, the diagonal Fisher information associated with
parameter $\theta_i$ is
\begin{equation}
    I(\theta_i)
    =
    \mathbb{E}
    \left[
        \left(
            \frac{\partial\log p(x;\boldsymbol{\theta})}
                 {\partial\theta_i}
        \right)^2
    \right].
    \label{eq:peft-fisher}
\end{equation}
An empirical approximation over $N$ observations is
\begin{equation}
    \widehat{I}(\theta_i)
    =
    \frac{1}{N}\sum_{j=1}^{N}
    \left(
        \frac{\partial\log p(x_j;\boldsymbol{\theta})}
             {\partial\theta_i}
    \right)^2.
    \label{eq:peft-empirical-fisher}
\end{equation}
Large diagonal Fisher values indicate that the log-likelihood is sensitive
to the corresponding parameters under the sampled distribution. The survey
discusses Fisher-based selection methods such as FishMask and Adafish, as
well as gradient-based selection methods \parencite{prottasha2025peft}.

\paragraph{Unstructured and structured selection.}
Unstructured selection chooses individual scalar parameters, potentially at
irregular positions throughout the model. BitFit, Child-Tuning, and other
sparsity-based methods illustrate this family. Structured selection chooses
coherent groups such as layers, heads, rows, columns, or blocks. A binary
mask can enforce a selection pattern:
\begin{equation}
    \boldsymbol{\theta}_{s}
    =
    \mathbf{M}\odot\boldsymbol{\theta},
    \qquad
    \mathbf{M}\in\{0,1\}^{P}.
    \label{eq:peft-structured-mask}
\end{equation}
Structured patterns can be easier to implement efficiently because they
preserve contiguous parameter groups. Unstructured sparsity may select more
precisely, but irregular memory access and indexing overhead can reduce the
hardware benefit. The survey discusses methods including UBitFit, LT-SFT,
Child-Tuning, RoCoFT, Xattn Tuning, and structured masking approaches
\parencite{prottasha2025peft}.

\paragraph{Low-rank reparameterization.}
Reparameterized PEFT constrains the weight update rather than selecting
individual original parameters. Consider a linear transformation
\begin{equation}
    \mathbf{y}
    =
    \mathbf{W}_0\mathbf{x},
    \qquad
    \mathbf{W}_0
    \in
    \mathbb{R}^{d_{\mathrm{out}}\times d_{\mathrm{in}}}.
    \label{eq:peft-linear-base}
\end{equation}
Full fine-tuning permits an arbitrary matrix update
$\Delta\mathbf{W}\in
\mathbb{R}^{d_{\mathrm{out}}\times d_{\mathrm{in}}}$, containing
$d_{\mathrm{out}}d_{\mathrm{in}}$ scalar degrees of freedom. Low-Rank
Adaptation (LoRA) parameterizes the update as
\begin{equation}
    \Delta\mathbf{W}
    =
    \mathbf{B}\mathbf{A},
    \qquad
    \mathbf{B}\in\mathbb{R}^{d_{\mathrm{out}}\times r},
    \quad
    \mathbf{A}\in\mathbb{R}^{r\times d_{\mathrm{in}}},
    \label{eq:peft-lora-factorization}
\end{equation}
where $r\ll\min(d_{\mathrm{in}},d_{\mathrm{out}})$. The effective weight is
\begin{equation}
    \mathbf{W}
    =
    \mathbf{W}_0+\mathbf{B}\mathbf{A}.
    \label{eq:peft-lora-weight}
\end{equation}
The number of trainable parameters is
\begin{equation}
    P_{\mathrm{LoRA}}
    =
    r(d_{\mathrm{out}}+d_{\mathrm{in}}),
    \label{eq:peft-lora-parameters}
\end{equation}
instead of $d_{\mathrm{out}}d_{\mathrm{in}}$. For a square $d\times d$
matrix, the ratio is
\begin{equation}
    \frac{P_{\mathrm{LoRA}}}{P_{\mathrm{FT}}}
    =
    \frac{2dr}{d^2}
    =
    \frac{2r}{d}.
    \label{eq:peft-lora-ratio}
\end{equation}
For example, with $d=4096$ and $r=8$, full fine-tuning of this matrix requires
$4096^2=16\,777\,216$ trainable values, whereas LoRA requires
$2(4096)(8)=65\,536$, approximately $0.39\%$ as many.

The matrix rank of a product cannot exceed the rank of either factor:
\begin{equation}
    \operatorname{rank}(\Delta\mathbf{W})
    =
    \operatorname{rank}(\mathbf{B}\mathbf{A})
    \leq r.
    \label{eq:peft-lora-rank}
\end{equation}
LoRA therefore restricts the update to a low-dimensional matrix subspace.
This reduces parameter count and can act as an implicit regularizer, but it
also limits the set of updates that can be represented. The rank $r$ controls
the capacity-efficiency trade-off. The survey describes an initialization in
which one factor is initialized to zero so that the initial low-rank update
is zero, preserving the pretrained transformation at the start of training
\parencite{prottasha2025peft}.

\paragraph{Dynamic rank and LoRA variants.}
A fixed rank allocates the same maximum update dimension throughout training.
Dynamic methods instead adapt rank or rank allocation according to training
signals or task-specific criteria. The survey discusses DyLoRA, AdaLoRA,
SLORA, CapaBoost, and AutoLoRA. A schematic representation is
\begin{equation}
    r_t=f(\|\nabla_t\|),
    \qquad t\in\{1,\ldots,T\},
    \label{eq:peft-dynamic-rank}
\end{equation}
where $f$ denotes a rank-allocation rule and $\nabla_t$ denotes gradient
information at iteration $t$. This equation expresses the general idea of
adaptive rank selection; it is not a universal update rule shared by every
dynamic-rank method.

The survey also reviews regularized and task-specific variants, including
LoRA Dropout, Laplace LoRA, Predict LoRA, LoRA++ and continual-learning
variants. Quantization can be combined with low-rank adaptation. QLoRA, for
example, keeps a quantized frozen base model and trains low-rank adaptation
parameters. SVDQUANT combines singular-value decomposition and quantization.
Variational LoRA introduces a probabilistic treatment of the adapted weights,
with the Bayesian relation
\begin{equation}
    p(\mathbf{W}\mid\mathcal{D})
    \propto
    p(\mathcal{D}\mid\mathbf{W})p(\mathbf{W}),
    \qquad
    \mathbf{W}=\mathbf{W}_0+\mathbf{B}\mathbf{A}.
    \label{eq:peft-variational-lora}
\end{equation}
These variants target different limitations, including rank allocation,
regularization, quantized storage, uncertainty, and continual adaptation
\parencite{prottasha2025peft}.

\paragraph{Hybrid, unified, and expert-based approaches.}
Hybrid methods combine adaptation mechanisms. UniPELT, for example, combines
LoRA, prefix tuning, and adapters within one framework. A general weighted
combination of $n$ adaptation contributions can be written as
\begin{equation}
    \Delta\boldsymbol{\theta}_{\mathrm{hybrid}}
    =
    \sum_{i=1}^{n}\alpha_i
    \Delta\boldsymbol{\theta}_i,
    \label{eq:peft-hybrid-update}
\end{equation}
where $\alpha_i$ controls the contribution of mechanism $i$. Unified
frameworks seek to coordinate or select among parameter-efficient modules
rather than committing to one fixed adaptation mechanism.

Mixture-of-experts-based PEFT extends this idea by maintaining several
specialized adaptation components. If expert $i$ supplies a low-rank update
$\mathbf{B}_i\mathbf{A}_i$, a routed update can be represented as
\begin{equation}
    \Delta\mathbf{W}(x)
    =
    \sum_{i=1}^{E}
    \alpha_i(x)\mathbf{B}_i\mathbf{A}_i,
    \qquad
    \sum_{i=1}^{E}\alpha_i(x)=1,
    \label{eq:peft-moe-lora}
\end{equation}
where the routing coefficients depend on the input or token representation.
Sparse routing can activate only a subset of experts. The survey discusses
MoE LoRA, MixLoRA, MoLoRA, MOA, and MoLE as examples of expert-based
adaptation \parencite{prottasha2025peft}.

\subsection{Sparse Matrix Tuning and Comparative Analysis}
\label{subsec:peft-smt-comparison}

Low-rank methods constrain the rank of the update matrix. Sparse Matrix
Tuning (SMT), introduced by \textcite{he2025smt}, takes a different approach:
it selects a small number of submatrices in the original weight matrices and
updates those blocks directly. The selected blocks need not be low-rank.
SMT therefore constrains the support of the update, rather than requiring the
entire update matrix to have low rank.

\paragraph{Task objective and block partitioning.}
Let $P_{\Phi_0}(y\mid x)$ be a pretrained autoregressive language model with
parameters $\Phi_0$, and let
$\mathcal{Z}=\{(x_i,y_i)\}_{i=1}^{N}$ be a downstream dataset. A conditional
language-modeling objective is
\begin{equation}
    \max_{\Theta}
    \sum_{(x,y)\in\mathcal{Z}}
    \sum_{t=1}^{|y|}
    \log
    P_{\Phi_0+\Delta\Phi(\Theta)}
    (y_t\mid x,y_{<t}),
    \label{eq:smt-objective}
\end{equation}
where $\Theta$ denotes the trainable adaptation variables. For LoRA,
$\Theta$ contains low-rank factors. For SMT, $\Theta$ contains the selected
submatrices.

Consider a pretrained weight matrix
$\mathbf{W}_0\in\mathbb{R}^{d\times k}$. Partition it into non-overlapping
blocks of size $l\times l$, assuming for the moment that $l$ divides both
dimensions. The number of blocks is
\begin{equation}
    N_{\mathrm{blocks}}
    =
    \frac{d}{l}\frac{k}{l}
    =
    \frac{dk}{l^2}.
    \label{eq:smt-number-blocks}
\end{equation}
If $\mathcal{S}$ is the selected set of $m$ blocks, the adapted weight is
\begin{equation}
    \mathbf{W}
    =
    \mathbf{W}_0+\Delta\mathbf{W}_{\mathcal{S}},
    \qquad
    \Delta\mathbf{W}_{\mathcal{S},ij}=0
    \quad\text{outside }\mathcal{S}.
    \label{eq:smt-sparse-update}
\end{equation}
The number of trainable values in this matrix is
\begin{equation}
    P_{\mathrm{SMT}}=ml^2,
    \qquad
    m\ll N_{\mathrm{blocks}}.
    \label{eq:smt-parameter-count}
\end{equation}
Within a selected block, the update can be a general $l\times l$ matrix.
Thus, a selected block can retain full local rank even though only a small
fraction of the complete matrix is trainable. Figure
\ref{fig:lora-versus-smt} contrasts the low-rank and submatrix update
principles.

\begin{figure}[htbp]
    \centering
    % Original conceptual comparison; uses the standard LaTeX picture environment.
\setlength{\unitlength}{1mm}
\begin{picture}(150,58)
  \put(18,52){\makebox(48,5){\textbf{LoRA: low-rank update}}}
  \put(92,52){\makebox(48,5){\textbf{SMT: selected blocks}}}

  % LoRA panel.
  \put(3,31){\framebox(22,11){\shortstack{Frozen\\$\mathbf{W}_0$}}}
  \put(31,42){\framebox(17,8){\scriptsize Trainable $\mathbf{A}$}}
  \put(31,25){\framebox(17,8){\scriptsize Trainable $\mathbf{B}$}}
  \put(48,46){\vector(1,0){7}}
  \put(48,29){\vector(1,0){7}}
  \put(55,34){\framebox(16,11){\shortstack{Low-rank\\$\mathbf{B}\mathbf{A}$}}}
  \put(25,36){\vector(1,0){6}}
  \put(71,39){\vector(1,0){4}}
  \put(3,10){\makebox(68,8){\shortstack{Update rank is at most $r$;\\factors add forward computation.}}}

  % SMT panel.
  \put(92,30){\framebox(36,16){}}
  \put(92,34){\line(1,0){36}}
  \put(92,38){\line(1,0){36}}
  \put(92,42){\line(1,0){36}}
  \put(101,30){\line(0,1){16}}
  \put(110,30){\line(0,1){16}}
  \put(119,30){\line(0,1){16}}
  \put(94,40){\makebox(6,3){\rule{3mm}{2mm}}}
  \put(112,36){\makebox(6,3){\rule{3mm}{2mm}}}
  \put(121,32){\makebox(6,3){\rule{3mm}{2mm}}}
  \put(92,47){\makebox(36,4){\scriptsize Partitioned base matrix $\mathbf{W}_0$}}
  \put(110,30){\vector(0,-1){6}}
  \put(88,14){\makebox(45,8){\shortstack{Selected blocks are trainable;\\unselected blocks remain frozen.}}}
\end{picture}

    \caption{Conceptual comparison of LoRA and Sparse Matrix Tuning. LoRA
    represents the update as a product of low-rank factors, whereas SMT
    partitions the original matrix and updates selected blocks. Original
    schematic based on \textcite{he2025smt}.}
    \label{fig:lora-versus-smt}
\end{figure}
% Alt text for accessibility register: The left panel shows a frozen base
% matrix combined with two trainable low-rank factors. The right panel shows
% a matrix divided into blocks, with only selected blocks updated.

\paragraph{Gradient-aware block selection.}
SMT must identify which blocks are most relevant to the downstream objective.
The paper compares activation-aware selection with gradient-aware selection.
In the reported experiments, gradient-aware selection is preferred. During
a warm-up phase, gradients from the relevant attention linear layers are
accumulated. Let $\mathbf{G}^{(t)}$ be the weight-gradient matrix at warm-up
iteration $t$. The accumulated gradient is
\begin{equation}
    \mathbf{G}^{\mathrm{acc}}
    =
    \sum_{t=1}^{T_w}\mathbf{G}^{(t)},
    \label{eq:smt-accumulated-gradient}
\end{equation}
where $T_w$ is the number of warm-up iterations. For block $B_j$, the score
used to rank blocks can be represented as
\begin{equation}
    s_j
    =
    \frac{1}{l^2}
    \sum_{(a,b)\in B_j}
    \left|G^{\mathrm{acc}}_{ab}\right|.
    \label{eq:smt-block-score}
\end{equation}
The absolute value prevents positive and negative entries from cancelling
when the magnitude of accumulated gradient information is assessed. The
blocks are ranked by $s_j$, and the highest-scoring blocks are retained until
the specified trainable-parameter budget is reached. The selection budget is
a method configuration; the exact trainable fraction varies across the
reported experiments.

The paper reports fewer than 100 warm-up iterations in its experiments. It
uses 100 warm-up iterations for the commonsense experiments and 25 for the
Math10K experiments. The block dimension is $256\times256$ in the reported
LLaMA experiments. This dimension is selected because it is a common factor
of the row and column dimensions of the relevant LLaMA linear layers, avoiding
remainder blocks. The authors also report that activation-aware selection
substantially underperforms gradient-aware selection in their ablation study
\parencite{he2025smt}.

\paragraph{Partial weight-gradient computation.}
The computational advantage of SMT depends on changing the backward
calculation, not merely freezing parameters after a conventional full
backward pass. Consider a linear layer
\begin{equation}
    \mathbf{Z}=\mathbf{W}\mathbf{X},
    \label{eq:smt-linear-layer}
\end{equation}
where $\mathbf{W}\in\mathbb{R}^{d_{\mathrm{out}}\times d_{\mathrm{in}}}$,
$\mathbf{X}\in\mathbb{R}^{d_{\mathrm{in}}\times n}$, and
$\mathbf{Z}\in\mathbb{R}^{d_{\mathrm{out}}\times n}$. Define
$\mathbf{G}=\partial\mathcal{L}/\partial\mathbf{Z}$. The standard matrix
derivatives are
\begin{align}
    \frac{\partial\mathcal{L}}{\partial\mathbf{W}}
    &=
    \mathbf{G}\mathbf{X}^{\mathsf T},
    \label{eq:smt-weight-gradient}\\
    \frac{\partial\mathcal{L}}{\partial\mathbf{X}}
    &=
    \mathbf{W}^{\mathsf T}\mathbf{G}.
    \label{eq:smt-input-gradient}
\end{align}
These derivatives have different roles. The first determines the weight
update; the second propagates the loss gradient to earlier layers. Freezing
a weight matrix removes the need to calculate its weight gradient, but the
input gradient may still be necessary if earlier components are trainable.

For a selected block with row indices $I$ and column indices $J$, the
required weight gradient is
\begin{equation}
    \frac{\partial\mathcal{L}}
         {\partial\mathbf{W}_{I,J}}
    =
    \mathbf{G}_{I,:}\mathbf{X}_{J,:}^{\mathsf T}.
    \label{eq:smt-selected-block-gradient}
\end{equation}
Only the corresponding rows of the output gradient and columns of the
activation are required for this block's weight-gradient calculation. SMT
therefore saves selected activation components and computes only the
selected weight gradients. It does not eliminate every backward operation:
the activation-gradient calculations that remain necessary for propagation
through the computational graph are retained.

The implementation replaces selected linear layers with customized sparse
linear layers. After selection, the chosen blocks are gathered and
concatenated into a dense tensor of shape $m\times l\times l$. The method
does not rely on general sparse matrix-matrix multiplication for these
updates; it uses dense matrix multiplication within the selected blocks.
This is a hardware-aware choice because irregular sparse indexing can
introduce overhead, whereas dense matrix multiplication is highly optimized
on contemporary accelerators. Figure \ref{fig:smt-workflow} summarizes the
selection and update process.

\begin{figure}[htbp]
    \centering
    % Original conceptual workflow; uses the standard LaTeX picture environment.
\setlength{\unitlength}{1mm}
\begin{picture}(150,58)
  \put(2,35){\framebox(25,13){\shortstack{\scriptsize Pretrained\\\scriptsize model + data}}}
  \put(31,35){\framebox(25,13){\shortstack{\scriptsize Warm-up\\\scriptsize backpropagation}}}
  \put(60,35){\framebox(25,13){\shortstack{\scriptsize Accumulate\\\scriptsize gradients}}}
  \put(89,35){\framebox(25,13){\shortstack{\scriptsize Score each\\\scriptsize submatrix}}}
  \put(118,35){\framebox(28,13){\shortstack{\scriptsize Select top\\\scriptsize blocks}}}
  \put(27,41){\vector(1,0){4}}
  \put(56,41){\vector(1,0){4}}
  \put(85,41){\vector(1,0){4}}
  \put(114,41){\vector(1,0){4}}

  \put(118,12){\framebox(28,13){\shortstack{\scriptsize Partial weight\\\scriptsize gradients}}}
  \put(89,12){\framebox(25,13){\shortstack{\scriptsize Update selected\\\scriptsize blocks only}}}
  \put(60,12){\framebox(25,13){\shortstack{\scriptsize Frozen base\\\scriptsize parameters}}}
  \put(132,35){\line(0,-1){10}}
  \put(132,25){\vector(0,-1){3}}
  \put(118,18){\vector(-1,0){4}}
  \put(89,18){\vector(-1,0){4}}
  \put(60,18){\vector(-1,0){4}}
  \put(2,4){\makebox(144,5){\scriptsize Warm-up selects the blocks; subsequent training avoids unselected weight gradients.}}
\end{picture}

    \caption{SMT workflow. A short warm-up phase accumulates gradients,
    submatrix scores determine the selected blocks, and a customized
    backward implementation computes and updates only the selected
    submatrix gradients. The diagram is an original schematic based on
    \textcite{he2025smt}.}
    \label{fig:smt-workflow}
\end{figure}
% Alt text for accessibility register: A flow diagram connects downstream
% data and a pretrained model to warm-up gradient accumulation, block
% scoring, top-block selection, and partial-gradient updates.

\paragraph{Why attention and value projections?}
The SMT experiments freeze MLP layers and apply their selected updates to
the query, key, and value projections in attention. The authors compare
parameter budgets allocated to MLPs and attention. On LLaMA-7B, with a
reported trainable budget of approximately $0.84\%$, allocating all selected
parameters to attention produces an average commonsense accuracy of
$78.7\%$, compared with $76.7\%$ when the same budget is allocated to MLPs.
Intermediate allocations produce intermediate results in the reported
ablation \parencite{he2025smt}.

The authors further compare the query, key, and value projections. Their
reported average accuracies are $73.9\%$ for key-only tuning, $75.3\%$ for
query-only tuning, $78.5\%$ for value-only tuning, and $78.7\%$ for joint
query-key-value selection. SMT assigns approximately $95.17\%$ of its
selected trainable parameters to value projections in the reported
visualization. These findings are empirical observations for the studied
LLaMA configurations and downstream benchmarks; they should not be
generalized into a universal claim that value projections are always the
most important parameters in every Transformer.

The attention computation helps explain why gradient magnitudes can differ
among these projections. For scaled dot-product attention,
\begin{equation}
    \mathbf{S}
    =
    \frac{\mathbf{Q}\mathbf{K}^{\mathsf T}}{\sqrt{d_k}},
    \qquad
    \mathbf{P}
    =
    \operatorname{softmax}(\mathbf{S}),
    \qquad
    \mathbf{O}=\mathbf{P}\mathbf{V}.
    \label{eq:smt-attention}
\end{equation}
For a row of softmax probabilities,
\begin{equation}
    \frac{\partial p_i}{\partial s_j}
    =
    p_i(\delta_{ij}-p_j).
    \label{eq:smt-softmax-jacobian}
\end{equation}
If one probability approaches one and the others approach zero, the
corresponding softmax derivatives approach zero. This is softmax saturation.
The gradient with respect to the value matrix follows directly through the
weighted aggregation:
\begin{equation}
    \frac{\partial\mathcal{L}}{\partial\mathbf{V}}
    =
    \mathbf{P}^{\mathsf T}
    \frac{\partial\mathcal{L}}{\partial\mathbf{O}}.
    \label{eq:smt-value-gradient}
\end{equation}
By contrast, gradients with respect to $\mathbf{Q}$ and $\mathbf{K}$ pass
through the softmax Jacobian. For one query vector $\mathbf{q}$ and key
vectors $\mathbf{k}_j$, the score derivative is
\begin{equation}
    \frac{\partial s_j}{\partial\mathbf{q}}
    =
    \frac{\mathbf{k}_j}{\sqrt{d_k}},
    \qquad
    \frac{\partial s_j}{\partial\mathbf{k}_j}
    =
    \frac{\mathbf{q}}{\sqrt{d_k}}.
    \label{eq:smt-qk-score-derivative}
\end{equation}
Consequently, saturation can attenuate the gradients reaching query and key
projections. This provides a mechanism consistent with the gradient
differences reported by the SMT authors, although the exact magnitude also
depends on the data, model state, and attention distribution. The scaling
factor $1/\sqrt{d_k}$ controls score variance under the usual independent,
zero-mean component assumptions. If the query and key components have
variances $\sigma_Q^2$ and $\sigma_K^2$, respectively, then
\begin{equation}
    \operatorname{Var}(\mathbf{q}^{\mathsf T}\mathbf{k})
    =
    d_k\sigma_Q^2\sigma_K^2,
    \qquad
    \operatorname{Var}
    \left(
        \frac{\mathbf{q}^{\mathsf T}\mathbf{k}}{\sqrt{d_k}}
    \right)
    =
    \sigma_Q^2\sigma_K^2.
    \label{eq:smt-attention-variance}
\end{equation}
The scaling controls the growth of score variance with $d_k$, but does not
guarantee that attention probabilities cannot become saturated.

\paragraph{Reported performance and system efficiency.}
The commonsense results reported by \textcite{he2025smt} show that SMT can
approach full fine-tuning while updating a small fraction of parameters.
For LLaMA-7B, the reported average accuracy is $81.4\%$ for full fine-tuning.
At approximately $0.84\%$ trainable parameters, SMT obtains $78.7\%$,
compared with $76.3\%$ for LoRA, $78.4\%$ for DoRA, and $76.5\%$ for SpIEL.
With a larger SMT budget of approximately $4.91\%$, the reported average
reaches $81.4\%$. These figures are benchmark-specific and should be read
together with the stated model, dataset, and hyperparameter configuration.

\begin{table}[htbp]
    \centering
    \caption{LLaMA-7B average accuracy on the commonsense reasoning
    benchmark reported by \textcite{he2025smt}. The trainable-parameter
    percentages and accuracy values are reproduced from the paper's
    comparison table.}
    \label{tab:smt-commonsense-results}
    \begin{tabularx}{\textwidth}{lcc}
        \toprule
        Method & Trainable parameters (\%) & Average accuracy (\%) \\
        \midrule
        Full fine-tuning & 100.00 & 81.4 \\
        LoRA & 0.83 & 76.3 \\
        DoRA & 0.84 & 78.4 \\
        SpIEL & 0.84 & 76.5 \\
        SMT & 0.84 & 78.7 \\
        SMT (larger budget) & 4.91 & 81.4 \\
        \bottomrule
    \end{tabularx}
\end{table}

For the LLaMA-7B timing comparison, the paper reports experiments on four
NVIDIA A100 40 GB GPUs. Under that configuration, full fine-tuning takes
$243.84$ seconds, while SMT with a $1.26\%$ trainable-parameter budget takes
$16.68$ seconds, corresponding to a reported $14.6\times$ speedup. LoRA,
DoRA, and SpIEL take $17.82$, $18.04$, and $25.45$ seconds, respectively,
under the reported comparison settings. These wall-clock values are
configuration-specific, not universal speedup guarantees. The paper also
reports a $67\%$ reduction in GPU memory footprint relative to full
fine-tuning and a $99.5\%$ reduction in optimizer memory under its stated
analysis \parencite{he2025smt}.

\paragraph{Complexity comparison.}
The parameter count is only one part of the cost of an adaptation method.
A useful comparison distinguishes trainable parameters, additional
parameters, forward-pass overhead, backward-pass work, and the storage
required by optimizer state. Table \ref{tab:peft-complexity} summarizes
representative parameter-count scalings. These expressions are schematic:
the exact count depends on how many layers and matrices are adapted, whether
biases are included, and the implementation details.

\begin{table}[htbp]
    \centering
    \caption{Representative trainable-parameter scalings for PEFT methods.
    Here $P$ is the base-model parameter count, $L$ is the number of adapted
    layers, $d$ is the hidden dimension, $d_b$ is the adapter bottleneck,
    $m$ is the prompt length, and $r$ is the LoRA rank.}
    \label{tab:peft-complexity}
    \begin{tabularx}{\textwidth}{lXX}
        \toprule
        Method & Approximate trainable parameters & Principal restriction \\
        \midrule
        Full fine-tuning & $P$ & No restriction on the update \\
        Serial adapters & $2Ldd_b$ & Bottleneck transformations \\
        Soft prompt tuning & $md$ & Input-level prompt vectors \\
        Prefix tuning & $\mathcal{O}(Lmd)$ & Layer-wise attention prefixes \\
        LoRA & $r(d_{\mathrm{in}}+d_{\mathrm{out}})$ per matrix & Low-rank update \\
        Selective tuning & $q\ll P$ & Selected original parameters \\
        SMT & $ml^2$ per matrix & Selected full-rank blocks \\
        Propulsion & $\mathcal{O}(d)$ & Feature-wise scaling \\
        \bottomrule
    \end{tabularx}
\end{table}

For LoRA and adapters, additional trainable modules are introduced and
executed in the forward path unless they are merged or otherwise transformed
for inference. Selective tuning and SMT update selected portions of the
original parameters. Prompt methods add trainable representations and may
increase the effective sequence length. Quantization-based methods address
the numerical precision and storage of the frozen base model, which is a
different axis of efficiency from reducing the number of trainable
parameters. The survey's complexity discussion distinguishes total
trainable parameters from additional parameters and emphasizes that
parameter count alone does not capture all deployment costs
\parencite{prottasha2025peft}.

\subsection{Limitations, Research Directions, and Connection to Scaling}
\label{subsec:peft-limitations-outlook}

PEFT trades unrestricted adaptation capacity for a smaller optimization
problem. This trade-off can improve efficiency and may reduce estimation
variance on low-resource tasks, but a restrictive parameterization can
increase approximation error when the downstream task requires substantial
changes to the pretrained representation. The survey notes that PEFT can
underperform when a task demands significant adaptation or differs sharply
from the pretraining application. Adapter dimensions, prompt lengths, rank
values, layer placements, and selection thresholds can also be sensitive
hyperparameters \parencite{prottasha2025peft}.

The statistical trade-off can be expressed conceptually as a balance between
approximation error and estimation error. A highly flexible full fine-tuning
procedure can represent a broad set of task-specific changes, but estimating
a large number of parameters from limited data can increase variance. A
constrained PEFT method reduces the dimension of the adaptation problem, but
may not contain an update that adequately represents the target task. The
best method is therefore not necessarily the one with the fewest trainable
parameters. It is the method that meets the task's performance requirement
within the available compute, memory, storage, and latency budgets.

Several research questions follow from this perspective. A theoretical
account of parameter influence could connect the mutual information between
adaptation modules and predictions with the curvature of the loss landscape.
For example, the Hessian
\begin{equation}
    \mathbf{H}
    =
    \nabla_{\boldsymbol{\theta}}^2
    \mathcal{L}_{\mathrm{task}}
    \label{eq:peft-hessian}
\end{equation}
describes local second-order curvature. Studying how low-dimensional
adaptation subspaces align with sensitive directions of this curvature could
help explain why some small updates are effective. The survey identifies
theoretical understanding of parameter influence as an open research
direction, rather than presenting a completed general theory
\parencite{prottasha2025peft}.

Layer-wise sensitivity is another open problem. If $q_l$ parameters are
assigned to layer $l$, a simplified allocation problem is
\begin{equation}
    \max_{\{q_l\}_{l=1}^{L}}
    \sum_{l=1}^{L}\Delta\mathcal{P}_l(q_l)
    \quad
    \text{subject to}
    \quad
    \sum_{l=1}^{L}q_l\leq Q,
    \label{eq:peft-layer-allocation}
\end{equation}
where $\Delta\mathcal{P}_l(q_l)$ denotes the performance contribution
associated with the parameter budget assigned to layer $l$, and $Q$ is the
total budget. This formulation makes explicit why uniform allocation need
not be optimal. Jacobian analysis and Fisher-information estimates are
possible tools for identifying sensitive layers. SMT provides one empirical
example of nonuniform allocation by selecting blocks according to their
gradient scores \parencite{prottasha2025peft,he2025smt}.

Task-aware and domain-specific adaptation also remains important. Legal,
medical, mathematical, and multilingual tasks can impose different
representational requirements. Future PEFT designs may use structure-aware
prompts, hierarchical adapters, or task-conditioned reparameterization.
Multimodal and non-Transformer architectures introduce further challenges:
an adaptation must preserve alignment between modalities while modifying
only selected components. The survey identifies vision-language, speech,
video, graph, and recurrent architectures as areas in which PEFT principles
can be developed beyond conventional text-only Transformers
\parencite{prottasha2025peft}.

Continual learning requires a model to acquire new tasks without losing
previous capabilities. Dynamically growing adapter banks, adapter
composition, sparse adapter stacking, and delta compression are possible
directions, but they introduce decisions about module selection, interference,
and storage. Interpretability is also unresolved: a small adapter or prompt
is not automatically understandable. Attribution methods and activation
visualization are needed to explain how task-specific modules affect model
predictions.

Privacy-preserving and federated PEFT may reduce communication by exchanging
small adaptation updates rather than complete model copies. However,
parameter efficiency alone does not provide a privacy guarantee; adaptation
parameters can still reveal information about training examples. Differential
privacy and secure aggregation must be treated as additional mechanisms.
Standardized benchmarks are likewise necessary so that methods are compared
under consistent models, datasets, trainable budgets, hardware, latency, and
memory measurements. Hardware-aware and energy-aware design is essential
because mathematical sparsity or a reduced parameter count does not
automatically translate into faster execution. Finally, meta-PEFT aims to
learn how to select adapter placements, prompt structures, rank values, or
other adaptation configurations across tasks, reducing manual
hyperparameter search \parencite{prottasha2025peft}.

The role of PEFT in the optimization narrative is now clear. Pretraining
constructs a reusable parameterized model; full fine-tuning adapts that model
by optimizing its complete parameter vector; and PEFT restricts or
reparameterizes the task-specific update to reduce the cost of adaptation.
The different methods make distinct compromises between flexibility,
regularization, modularity, optimizer memory, and hardware efficiency.
SMT illustrates that the organization of the selected parameters and the
implementation of their gradients can matter as much as the abstract
parameter count.

This completes the discussion of optimization-side adaptation mechanisms.
The next stage of the book moves from the question of how a model is
optimized or adapted to the broader question of how model performance,
data, and computational resources relate as model scale changes. That
question leads naturally to the study of scaling laws.
